@article{iyengar2019origins,
  title={The origins and consequences of affective polarization in the United States},
  author={Iyengar, Shanto and Lelkes, Yphtach and Levendusky, Matthew and Malhotra, Neil and Westwood, Sean J},
  journal={Annual review of political science},
  volume={22},
  number={1},
  pages={129--146},
  year={2019},
  publisher={Annual Reviews}
}
Iyengar, Lelkes, Levendusky, Malhotra, and Westwood (2019), “The Origins and Consequences of Affective Polarization in the United States,” Annual Review of Political Science.
This is a standard review article. It states that ordinary Americans have “increasingly” come to dislike and distrust the opposing party, making it a very clean citation for the rise in affective polarization in the U.S.

@article{boxell2024cross,
  title={Cross-country trends in affective polarization},
  author={Boxell, Levi and Gentzkow, Matthew and Shapiro, Jesse M},
  journal={Review of Economics and Statistics},
  volume={106},
  number={2},
  pages={557--565},
  year={2024},
  publisher={MIT Press One Rogers Street, Cambridge, MA 02142-1209, USA journals-info~…}
}
Using data from twelve OECD countries over several decades, the authors find that the United States experienced the largest increase in affective polarization among the countries they study.

@article{bertrand2023coming,
  title={Coming apart? Cultural distances in the United States over time},
  author={Bertrand, Marianne and Kamenica, Emir},
  journal={American Economic Journal: Applied Economics},
  volume={15},
  number={4},
  pages={100--141},
  year={2023},
  publisher={American Economic Association 2014 Broadway, Suite 305, Nashville, TN 37203-2425}
}
This gives you an economics-style complement to the political science literature. It finds that differences in social attitudes by political ideology have increased over the last four decades. That is not identical to affective polarization, but it is very helpful if you want to support a broader claim about rising ideological or cultural polarization.


Pew Research Center (2023), “Americans’ Dismal Views of the Nation’s Politics.”
This is not a journal article, but it is very useful for the “through 2023” part. Pew reports that partisan divisions on issues are wider than a few decades ago and that many Americans hold deeply negative views of the other side. That gives you contemporary support that polarization remained high in 2023.

